%----------------------------------------------------------------------------------------
%	RESEARCH
%----------------------------------------------------------------------------------------

\section{Research Experience}

\researchentry
  {Water-Responsive Metal--Peptide Frameworks}
  {The University of North Carolina at Charlotte}
  {Volunteer Researcher}
  {Qiang Zhu}
  {Oct 2026 -- Present}
  {}
\begin{itemize}[
  leftmargin=1em,
  label={--},
  itemsep=2pt,
  topsep=0pt,
  parsep=0pt,
  partopsep=0pt
]
  \item Scoping a first reproducible simulation target, centered on the
  Zn--Gly--Thr versus Zn--Gly--Ser framework comparison (Yuan et al., 2025):
  identifying the input structures, interaction parameters, and experimental
  data required.

  \item In weekly meetings with graduate student Z.\,H.~Fahim, reviewing the
  literature on water-responsive peptide frameworks and MOFs, separating
  measured results from model-based explanations, and linking planned
  molecular dynamics and grand-canonical Monte Carlo calculations (including
  force-field fitting to DFT reference energies) to experiments that could
  validate them.
\end{itemize}

\vspace{0.2em}

\researchentry
  {RSC: Resolved Stellar Clusters}
  {American Museum of Natural History}
  {Volunteer Researcher}
  {Dave Zurek}
  {Sep 2026 -- Present}
  {}
\begin{itemize}[
  leftmargin=1em,
  label={--},
  itemsep=2pt,
  topsep=0pt,
  parsep=0pt,
  partopsep=0pt
]
  \item Audited and corrected errors in the statistical calculations of a
  Python code that assigns stellar-cluster membership probabilities from Gaia
  astrometry, using an expectation--maximization mixture model with per-star
  measurement errors.
\end{itemize}

\vspace{0.2em}

\researchentry
  {Starspot Inference Using Light Curve Inversion Techniques}
  {American Museum of Natural History}
  {M.S. Thesis Researcher}
  {Lucy Lu}
  {Aug 2023 -- Sep 2025}
  {\thesislink{https://academicworks.cuny.edu/gc_etds/6480/}}
\begin{itemize}[
  leftmargin=1em,
  label={--},
  itemsep=2pt,
  topsep=0pt,
  parsep=0pt,
  partopsep=0pt
]
  \item Built Python workflows with \texttt{fleck} and \texttt{starry} to
  compare forward and inverse models across 1,000 synthetic, evolving-starspot
  light curves generated with \texttt{butterpy}.

  \item Tested recovery of rotation, activity indicators, and surface structure
  across changes in photometric noise, viewing inclination, and spot
  configuration.

  \item Found that rotation and statistical activity indicators were
  recoverable in the tested cases, while individual spot locations and surface
  maps remained non-unique.
\end{itemize}

\vspace{0.2em}

\researchentry
  {Towards Mapping Brown Dwarf and Giant Exoplanet Atmospheres}
  {American Museum of Natural History}
  {BDNYC Researcher}
  {Johanna Vos}
  {May 2021 -- Aug 2023}
  {\githubrepo{https://github.com/MohammadRefat23/StarryAtmospheres}}
\begin{itemize}[
  leftmargin=1em,
  label={--},
  itemsep=2pt,
  topsep=0pt,
  parsep=0pt,
  partopsep=0pt
]
  \item Fit light curves from three-dimensional atmospheric simulations and
  20-hour Spitzer observations of two brown dwarfs with \texttt{starry}
  spherical-harmonic surface maps.

  \item Found degree $l=5$ sufficient to reproduce the relevant light-curve
  structure and identified evolving clouds as a limitation of static-map fits,
  motivating rotation-by-rotation mapping.
\end{itemize}

\vspace{0.2em}

\researchentry
  {Chemodynamically Characterizing the Jhelum Stellar Stream with APOGEE-2}
  {Sloan Digital Sky Survey Faculty and Student Team}
  {SDSS FAST Student Researcher}
  {Allyson Sheffield}
  {May 2020 -- May 2021}
  {%
    \githubrepo{https://github.com/MohammadRefat23/jhelum-stellar-stream}
    \hspace{0.15em}
    \paperlink{https://doi.org/10.3847/1538-4357/abee93}
  }
\begin{itemize}[
  leftmargin=1em,
  label={--},
  itemsep=2pt,
  topsep=0pt,
  parsep=0pt,
  partopsep=0pt
]
  \item Performed the chemical-abundance analysis for APOGEE-2 spectra of stars
  in six Jhelum fields, comparing [Ca/Fe], [C/Fe], [Na/Fe], [Mg/Fe], and
  [Ni/Fe] against [Fe/H] with published high-resolution Jhelum members to
  place the stream among accreted halo populations.

  \item Contributed to the collaboration's identification of a new metal-poor
  red-giant member ([Fe/H]~$\approx -2.2$) whose abundances favored an accreted
  dwarf-galaxy origin, with a globular-cluster origin not excluded.

  \item In 2026, rebuilt the workflow as a documented, tested Python package
  using public SDSS DR17 and VizieR data.
\end{itemize}

\vspace{0.2em}

\researchentry
  {Where Were Stars Born in the Milky Way?}
  {American Museum of Natural History}
  {AstroCom NYC Scholar}
  {Allyson Sheffield}
  {Aug 2018 -- May 2020}
  {}
\begin{itemize}[
  leftmargin=1em,
  label={--},
  itemsep=2pt,
  topsep=0pt,
  parsep=0pt,
  partopsep=0pt
]
  \item Processed and normalized high-resolution echelle spectra with
  \texttt{SMHr}, measured absorption-line equivalent widths, and iteratively
  fit stellar parameters and elemental abundances for chemical-tagging
  analysis.

  \item Validated the workflow against published Arcturus measurements,
  recovering effective temperature within 50\,K and metallicity within
  0.07\,dex of reference values.
\end{itemize}

\vspace{0.2em}

\researchentry
  {Two-Point Statistics in the Star-Forming Interstellar Medium}
  {Flatiron Institute, Center for Computational Astrophysics}
  {AstroCom NYC Scholar}
  {Chang-Goo Kim}
  {Jun 2018 -- Aug 2018}
  {\githubrepo{https://github.com/MohammadRefat23/tigress-metallicity-correlations}}
\begin{itemize}[
  leftmargin=1em,
  label={--},
  itemsep=2pt,
  topsep=0pt,
  parsep=0pt,
  partopsep=0pt
]
  \item Developed Python routines to compute and visualize two-point
  metallicity correlations in TIGRESS magnetohydrodynamic simulations of the
  star-forming interstellar medium.

  \item Measured a characteristic metallicity half-correlation length of
  $120\pm25\,\mathrm{pc}$ and estimated a turbulent metal-mixing timescale of
  $\sim10\,\mathrm{Myr}$.
\end{itemize}

\vspace{0.2em}

\researchentry
  {Staying in Science}
  {American Museum of Natural History}
  {Research Contributor, NSF-Funded Study}
  {Preeti Gupta}
  {Aug 2017 -- Sep 2019}
  {}
\begin{itemize}[
  leftmargin=1em,
  label={--},
  itemsep=2pt,
  topsep=0pt,
  parsep=0pt,
  partopsep=0pt
]
  \item Wrote reflective research memos on survey responses from 470 alumni of
  NYC science research mentoring programs, examining which science practices
  and interactions alumni reported in college, internships, and work; the
  memos served as a validity and reliability check on the research team's
  analysis.

  \item Also took part as a study participant, contributing data for
  social-network maps of students' STEM education and career pathways.
\end{itemize}
